\documentclass[a4paper,12pt]{article}
\usepackage[utf8]{inputenc}
\usepackage[T1]{fontenc}
\usepackage[ngerman]{babel}
\usepackage{geometry}
\usepackage{booktabs}
\usepackage{amsmath}
\geometry{margin=2.5cm}

\begin{document}

\section*{Geometry Aware Physics Informed Neural Network Surrogate for Solving
Navier--Stokes Equation (GAPINN)}

\textbf{Autoren:} Jan Oldenburg, Finja Borowski, Alper \"Oner, Klaus-Peter Schmitz,
Michael Stiehm \\
\textbf{Journal:} Advanced Modeling and Simulation in Engineering Sciences \\
\textbf{Jahr:} 2022 \quad DOI: \texttt{10.1186/s40323-022-00221-z}

% -------------------------------------------------------------------
\subsection*{Motivation und Problemstellung}

Die Simulation von Str\"omungsph\"anomenen in komplexen, irregul\"aren Geometrien
stellt ein zentrales Problem in vielen Ingenieur- und Naturwissenschaften dar.
Klassische numerische L\"oser wie die Finite-Volumen-Methode (FVM), implementiert
etwa in OpenFOAM, liefern zwar pr\"azise Ergebnisse, erfordern jedoch f\"ur jede
neue Geometrie eine aufwendige Vernetzung (Meshing) sowie eine separate,
rechenintensive Simulation. Dies macht sie f\"ur zeitkritische Anwendungen --
etwa die intraoperative Str\"omungssimulation in der Kardiologie -- praktisch
ungeeignet.

Surrogate-Modelle bieten hier einen vielversprechenden Ausweg: Ein einmal
trainiertes Modell kann f\"ur neue Eingaben innerhalb von Millisekunden eine
N\"aherungsl\"osung liefern. W\"ahrend datengetriebene Surrogate-Modelle
(z.\,B.\ auf Basis von U-Nets oder Fourier Neural Operators) gro\ss{}e
Trainingsmengen an Simulationsdaten ben\"otigen, verfolgen Physics-Informed
Neural Networks (PINNs) einen anderen Ansatz: Sie ersetzen oder erg\"anzen
Trainingsdaten durch die direkten Residuen der zugrundeliegenden
partiellen Differentialgleichungen (PDEs) in der Verlustfunktion.

Eine wesentliche offene Frage war bisher, ob PINNs auch f\"ur \emph{nicht
parametrisierte, irregul\"are Geometrien} trainiert werden k\"onnen -- also
f\"ur Formen, die sich nicht durch wenige Parameter (z.\,B.\ Radius einer
Engstelle) beschreiben lassen. Genau diese L\"ucke adressiert GAPINN.

% -------------------------------------------------------------------
\subsection*{Methodik: Das GAPINN-Framework}

GAPINN besteht aus drei separat trainierten Teilnetzen, die sequenziell
zusammenarbeiten:

\paragraph{1. Shape Encoding Network (SEN).}
Das SEN komprimiert die hochdimensionale Geometrieinformation -- repr\"asentiert
als Punktwolke der Randpunkte $\mathbf{X}_{j,b}$ -- in einen niedrigdimensionalen
Latentvektor $\mathbf{k} \in \mathbb{R}^{n_k}$ (hier: $n_k = 60$). Als
Architektur wird ein Variational Autoencoder (VAE) mit einer PointNet-artigen
Encoder-Struktur eingesetzt: vier 1D-Faltungsschichten (128, 128, 256, 512
Kan\"ale), gefolgt von einem kanalweisen Max-Pooling und zwei vollverkn\"upften
Netzwerken, die Mittelwert und Varianz der Posteriorverteilung sch\"atzen.
Die Verwendung eines VAE statt eines gew\"ohnlichen Autoencoders sorgt daf\"ur,
dass geometrisch \"ahnliche Formen auch im Latentraum nah beieinander liegen
(Regularisierung via KL-Divergenz), was die Interpolationsf\"ahigkeit verbessert.
Die PointNet-Architektur gew\"ahrleistet Permutationsinvarianz bez\"uglich der
Reihenfolge der Eingabepunkte. Nach dem Training des SEN werden dessen Gewichte
eingefroren.

\paragraph{2. Boundary Constraining Network (BCN).}
Das BCN sch\"atzt f\"ur jeden Raumpunkt $\mathbf{X}_{j,(i,b)}$ den minimalen
euklidischen Abstand zur n\"achsten Wand mit Haftbedingung ($U_\text{wall} = 0$).
Dieser Abstand dient als stetig differenzierbare ``Abstandsfunktion'', die im
dritten Netz genutzt wird, um Randbedingungen \emph{hart} zu erzwingen -- im
Gegensatz zu weichen Randbedingungen, die lediglich als zus\"atzlicher
Strafterm in die Verlustfunktion eingehen. Der entscheidende Vorteil: Da das BCN
differenzierbar ist, kann w\"ahrend des PINN-Trainings automatisches
Differenzieren durch die Randbedingungsenforcierung hindurch angewendet werden.

\paragraph{3. Physics Informed Neural Network (PINN).}
Das PINN ist ein vollverkn\"upftes Feedforward-Netzwerk, das als Eingabe die
Raumkoordinaten $\mathbf{X}_{j,(i,b)}$ und den Latentvektor $\mathbf{k}$
erh\"alt. Es gibt Geschwindigkeit $\mathbf{U}$ und Druck $p$ aus. Die
Verlustfunktion enth\"alt \emph{ausschlie\ss{}lich} die Residuen der
station\"aren, inkompressiblen Navier-Stokes-Gleichungen:
\begin{equation}
  \mathcal{L}_\text{phy}(\mathbf{W}, \mathbf{b}) =
  \bigl\|\nabla \mathbf{U}\bigr\|^2 +
  \Bigl\|(\mathbf{U} \cdot \nabla)\mathbf{U}
  + \tfrac{1}{\rho}\nabla p - \nu\nabla^2 \mathbf{U}\Bigr\|^2,
\end{equation}
berechnet mittels automatischer Differentiation. Es werden \emph{keine}
Referenzsimulationen als Trainingsziel verwendet (datenfreies Training).

\paragraph{Randbedingungsenforcierung.}
Die Ausgaben des PINN werden durch die harten Randbedingungsfunktionen des BCN
transformiert, sodass der endg\"ultige Output $\mathbf{U}_{j,(i,b)}$ die
Haftbedingung an W\"anden und das parabolische Einstr\"omprofil am Einlass
exakt erf\"ullt. Die Autoren zeigen, dass diese harte Enforcierung gegen\"uber
dem weichen Strafterm zu deutlich schnellerer Konvergenz und niedrigeren
Verlusten f\"uhrt.

\paragraph{Experimente und Ergebnisse.}
Das Framework wird auf 1000 synthetisch erzeugte 2D-Gef\"a\ss{}geometrien mit
Stenosen und Aneurysmen trainiert ($Re = 500$, station\"ar, laminar).
Validierung erfolgt an Geometrien, die nicht im Training enthalten waren.
Der Vergleich mit CFD-Simulationen (OpenFOAM, ANSYS Fluent) zeigt gute
qualitative \"Ubereinstimmung: MSE$_u = 3{,}61 \times 10^{-3}$,
MSE$_v = 2{,}30 \times 10^{-4}$, MSE$_p = 1{,}6 \times 10^{-2}$.
Die Inferenz f\"ur eine neue Geometrie dauert $0{,}32$ GPU-Sekunden,
verglichen mit 210 CPU-Sekunden f\"ur eine einzelne CFD-Simulation.
Ein Vanilla-PINN (ohne SEN/BCN) scheitert auf dieser Aufgabe vollst\"andig.

% -------------------------------------------------------------------
\subsection*{Bezug zur Masterarbeit}

GAPINN und die vorliegende Masterarbeit teilen dieselbe grundlegende Motivation:
Rechenintensive CFD-Simulationen f\"ur physikalische Systeme mit
\emph{wechselnden Geometrien} durch schnelle, neuronale Surrogate zu ersetzen.
Dennoch unterscheiden sich beide Arbeiten in Ansatz, Architektur und
Problemstellung erheblich. Die folgende Gegen\"uberstellung verdeutlicht
Gemeinsamkeiten und Unterschiede:

\begin{center}
\begin{tabular}{lll}
\toprule
\textbf{Aspekt} & \textbf{GAPINN} & \textbf{Masterarbeit} \\
\midrule
Physikalisches Regime & Navier-Stokes ($Re = 500$) & Stokes-Str\"omung ($Re \ll 1$) \\
Geometrievariabilit\"at & Unterschiedliche Kanalgeometrien & $1$ bis $N$ Kristalle \\
Architekturtyp & PINN (kein Trainingsdatensatz) & U-Net, FNO (datengetrieben) \\
Trainingsstrategie & Datenfreies, physik-getriebenes Training & LaMEM-Simulationsdaten \\
Geometrierepr\"asentation & Latentvektor via VAE/SEN & Bin\"armaske / SDF auf Grid \\
Ausgangs\-gr\"o\ss{}e & Geschwindigkeit + Druck & Stromfunktion $\psi$ \\
Randbedingungen & Hart erzwungen via BCN & Soft (Verlustterm) \\
\bottomrule
\end{tabular}
\end{center}

\medskip

Konkret lassen sich folgende inhaltliche Bez\"uge identifizieren:

\begin{enumerate}

\item \textbf{Geometrierepresentation als Schl\"usselherausforderung.}
GAPINN l\"ost das Problem der variablen Geometrie durch einen gelernten
Latentvektor (VAE/SEN), der die Kanalform in $\mathbb{R}^{60}$ komprimiert.
In der Masterarbeit wird ein komplementarer, gitterbasierter Ansatz verfolgt:
Die Kristallpositionen werden als Bin\"armaske oder Signed Distance Function
(SDF) auf dem $256 \times 256$-Grid kodiert und direkt als Eingabekanal
in das Netzwerk eingespeist. Beide Strategien adressieren dasselbe
grundlegende Problem -- wie vermittelt man dem Netzwerk effizient die
raumabh\"angige geometrische Struktur? --, verfolgen jedoch fundamental
verschiedene Wege: GAPINN abstrahiert die Geometrie global in einen
Vektor, w\"ahrend U-Net und FNO die geometrische Information lokal bzw.\
spektral auf dem vollen Grid verarbeiten.

\item \textbf{Generalisierung auf ungesehene Konfigurationen.}
GAPINN demonstriert Generalisierung auf neue Kanalgeometrien (Validierungsset),
die nicht im Training enthalten waren, sofern diese geometrisch \"ahnlich zu
Trainingsgeometrien sind. Die Masterarbeit stellt eine anspruchsvollere
Variante dieser Generalisierungsaufgabe: Die Modelle sollen nicht nur auf
neue \emph{Anordnungen} einer fixen Anzahl von Kristallen generalisieren,
sondern auf eine \emph{variable Anzahl} von Kristallen -- von 1 bis $N+m$.
Dies stellt h\"ohere Anforderungen an die Generalit\"at der gelernten
Repr\"asentation, da sich die Komplexit\"at des Str\"omungsfelds mit der
Kristallzahl qualitativ ver\"andert.

\item \textbf{Physikalische Konsistenz als Trainingsziel.}
GAPINN erzwingt physikalische Konsistenz durch die explizite Minimierung
der PDE-Residuen und die harte Randbedingungsenforcierung. In der Masterarbeit
wird ein alternativer, aber verwandter Weg beschritten: Die Vorhersage der
Stromfunktion $\psi$ gew\"ahrleistet die Inkompressibilit\"at
($\nabla \cdot \mathbf{u} = 0$) \emph{analytisch durch die mathematische
Konstruktion} des Ausgabefelds, da $\mathbf{u} = \nabla \times \psi$.
Dies entspricht dem Konzept der harten physikalischen Constraint-Enforcierung
bei GAPINN, wird aber durch die Wahl der Ausgabegr\"o\ss{}e statt durch ein
zus\"atzliches Netzwerk realisiert.

\item \textbf{Grenzen des PINN-Ansatzes und Motivation f\"ur datengetriebene
Modelle.}
GAPINN zeigt, dass das datenfreie Training auf Basis der PDE-Residuen auch
bei hoher geometrischer Variabilit\"at prinzipiell funktioniert, jedoch
sehr langen Trainingszeiten (ca.\ 13 h f\"ur 1000 Geometrien auf einer GPU)
erfordert und bisher nur f\"ur relativ einfache laminare Regime validiert
wurde. F\"ur die Masterarbeit -- in der Dutzende von Architekturen, Eingabe-
repr\"asentationen und Trainingsstrategien systematisch verglichen werden
sollen -- ist ein datengetriebener Ansatz praktikabler, da LaMEM-Simulationen
parallelisierbar erzeugt werden k\"onnen und das Training der Surrogate
wesentlich effizienter konvergiert. Dies zeigt, dass GAPINN und der Ansatz
der Masterarbeit als komplement\"are Strategien f\"ur \"ahnliche Probleme
zu verstehen sind.

\item \textbf{Relevanz f\"ur die Netzwerkarchitektur-Diskussion.}
GAPINN basiert auf vollverkn\"upften Netzwerken (FCNN) und ist damit
grunds\"atzlich verschieden von den konvolutionalen U-Net- und
spektralen FNO-Architekturen der Masterarbeit. Die Autoren selbst
diskutieren, dass f\"ur gitterbasierte Repr\"asentationen -- wie sie in
der Masterarbeit verwendet werden -- 2D-CNNs eine nat\"urliche Alternative
w\"aren. Dies best\"atigt retrospektiv die Architekturwahl der Masterarbeit:
U-Net und FNO k\"onnen geometrische Information effizienter verarbeiten,
wenn sie bereits auf einem strukturierten Gitter kodiert vorliegt.

\end{enumerate}

\subsection*{Einordnung in die Literatur}

GAPINN ist in Kapitel~2 der Masterarbeit im Abschnitt zu geometrie-bewussten
PINNs zu zitieren (zusammen mit Sun et al.\ 2020, PhyGeoNet von Gao et al.\
2021). Es belegt, dass das Problem der Generalisierung \"uber variable
Geometrien aktiv erforscht wird und dass verschiedene Strategien zur
Geometrierepr\"asentation -- von Latentvektoren bis zu gitterbasierten Masken
-- prinzipiell geeignet sind. Gegen\"uber der Masterarbeit fehlt GAPINN
jedoch die Untersuchung der \emph{skalaren Anzahlvariabilit\"at} von Objekten
sowie ein direkter Vergleich konkurrierender Netzwerkparadigmen (U-Net vs.\
FNO), was die Originalit\"at der Masterarbeit unterstreicht.

\end{document}